\begin{abstract}
Network slicing promises logically isolated, service-tailored virtual
networks on shared 5G infrastructure, but its operation hinges on a
question that is usually answered crudely: how much capacity should a
slice reserve for the next interval when future demand is uncertain?
Over-reservation wastes spectrum that other tenants could monetize;
under-reservation violates service-level agreements and degrades
user experience. Existing approaches occupy two unattractive extremes:
static margins sized from peak historical load, which waste 30\% or
more of capacity, and end-to-end learned controllers, whose policies
are opaque, hard to audit, and brittle under distribution shift. This
paper presents UCRA (Uncertainty-to-Reservation Closed-Loop Resource
Allocation), a five-stage framework that turns calibrated probabilistic
demand forecasts into capacity reservations through an explicit,
auditable transform $\Rt = \qtau + \kappa\,\spread\,(1 + w\,\rhot)$,
where $\qtau$ is a deep quantile forecast, $\spread$ its dispersion,
$\rhot$ an online risk indicator, and $\kappa$ a single operator-facing
risk dial. A self-evolving stage detects drift from reservation
feedback and fine-tunes the forecaster on a reservoir replay buffer
under a formally causal information-flow protocol that we prove and
enforce in tests: the model can only learn from labels that would
already exist in a live deployment. On three weeks of live 5G/4G/2G
radio access network performance counters (302{,}052 records, 75
sectors), UCRA holds \textbf{0.0\% reservation violations at 72.1\%
utilization} and 10.7\% waste, where a peak-sized static reservation
wastes 30.5\%, a median-forecast policy violates 48.0\% of intervals,
and even a cheating oracle quantile still violates 16.3\%. A
three-seed study confirms stability (0.0\% violations in every seed;
utilization $67.6\%\pm3.4$). Under an injected $+35\%$ demand surge,
the causal self-evolving loop cuts post-surge violations by half
($34.3\%\pm5.9 \rightarrow 17.2\%\pm4.9$ across seeds) relative to a
frozen forecaster, while firing zero spurious updates on clean data. On
an independent CTTC 5G slicing workload (229 steady-state snapshots),
the same transform reduces slice reservation violations to 1.2--2.7\%
across URLLC, eMBB, and mMTC where the operator's default sizing
violates 5.8--43.7\%. The full pipeline --- data ingestion, training,
closed-loop evaluation, leakage audit, and three-seed verification ---
is open and reproducible from the accompanying repository.
\end{abstract}

\begin{IEEEkeywords}
5G network slicing, resource reservation, quantile regression, deep
learning, uncertainty quantification, online learning, concept drift,
replay buffers, risk management, reproducibility.
\end{IEEEkeywords}
